\documentclass[aspectratio=169]{beamer}

\usepackage{iftex}
\providecommand{\dunenoirsetupfonts}{}
\newcommand{\DuneNoirLocalFontPath}{fonts}

\newcommand{\DuneNoirMonoFallback}{%
  \IfFontExistsTF{JetBrainsMono Nerd Font Propo}{
    \setmonofont{JetBrainsMono Nerd Font Propo}[Scale=0.92]
  }{
    \IfFontExistsTF{JetBrainsMono Nerd Font}{
      \setmonofont{JetBrainsMono Nerd Font}[Scale=0.92]
    }{
      \IfFontExistsTF{Liberation Mono}{
        \setmonofont{Liberation Mono}[Scale=0.92]
      }{
        \IfFontExistsTF{DejaVu Sans Mono}{
          \setmonofont{DejaVu Sans Mono}[Scale=0.92]
        }{
          \setmonofont{Courier New}[Scale=0.92]
        }
      }
    }
  }
}

\newcommand{\DuneNoirSansFallback}{%
  \IfFontExistsTF{IBM Plex Sans}{
    \setsansfont{IBM Plex Sans}
    \setmainfont{IBM Plex Sans}
  }{
    \IfFontExistsTF{Helvetica Neue}{
      \setsansfont{Helvetica Neue}
      \setmainfont{Helvetica Neue}
    }{
      \IfFontExistsTF{Liberation Sans}{
        \setsansfont{Liberation Sans}
        \setmainfont{Liberation Sans}
      }{
        \IfFontExistsTF{DejaVu Sans}{
          \setsansfont{DejaVu Sans}
          \setmainfont{DejaVu Sans}
        }{
          \setsansfont{Arial}
          \setmainfont{Arial}
        }
      }
    }
  }
}

\newcommand{\DuneNoirUseHelvetica}{%
  \renewcommand{\dunenoirsetupfonts}{%
    \IfFontExistsTF{Helvetica Neue}{
      \setsansfont{Helvetica Neue}
      \setmainfont{Helvetica Neue}
    }{
      \IfFontExistsTF{Helvetica}{
        \setsansfont{Helvetica}
        \setmainfont{Helvetica}
      }{
        \DuneNoirSansFallback
      }
    }
    \DuneNoirMonoFallback
  }%
}

\newcommand{\DuneNoirUseOpenSans}{%
  \renewcommand{\dunenoirsetupfonts}{%
    \IfFileExists{\DuneNoirLocalFontPath/OpenSans.ttf}{
      \setsansfont{Open Sans Custom}[
        Path=\DuneNoirLocalFontPath/,
        UprightFont={OpenSans.ttf},
        ItalicFont={OpenSans-Italic.ttf},
        BoldFont={OpenSans.ttf},
        BoldItalicFont={OpenSans-Italic.ttf},
        BoldFeatures={FakeBold=1.45},
        BoldItalicFeatures={FakeBold=1.45}
      ]
      \setmainfont{Open Sans Custom}[
        Path=\DuneNoirLocalFontPath/,
        UprightFont={OpenSans.ttf},
        ItalicFont={OpenSans-Italic.ttf},
        BoldFont={OpenSans.ttf},
        BoldItalicFont={OpenSans-Italic.ttf},
        BoldFeatures={FakeBold=1.45},
        BoldItalicFeatures={FakeBold=1.45}
      ]
    }{
      \DuneNoirSansFallback
    }
    \DuneNoirMonoFallback
  }%
}

\newcommand{\DuneNoirUseDosis}{%
  \renewcommand{\dunenoirsetupfonts}{%
    \IfFileExists{\DuneNoirLocalFontPath/Dosis.ttf}{
      \setsansfont{Dosis Custom}[
        Path=\DuneNoirLocalFontPath/,
        UprightFont={Dosis.ttf},
        ItalicFont={Dosis.ttf},
        BoldFont={Dosis.ttf},
        BoldItalicFont={Dosis.ttf},
        ItalicFeatures={FakeSlant=0.16},
        BoldFeatures={FakeBold=1.6},
        BoldItalicFeatures={FakeBold=1.6,FakeSlant=0.16}
      ]
      \setmainfont{Dosis Custom}[
        Path=\DuneNoirLocalFontPath/,
        UprightFont={Dosis.ttf},
        ItalicFont={Dosis.ttf},
        BoldFont={Dosis.ttf},
        BoldItalicFont={Dosis.ttf},
        ItalicFeatures={FakeSlant=0.16},
        BoldFeatures={FakeBold=1.6},
        BoldItalicFeatures={FakeBold=1.6,FakeSlant=0.16}
      ]
    }{
      \DuneNoirSansFallback
    }
    \DuneNoirMonoFallback
  }%
}

\newcommand{\DuneNoirUseKoHo}{%
  \renewcommand{\dunenoirsetupfonts}{%
    \IfFileExists{\DuneNoirLocalFontPath/KoHo-Regular.ttf}{
      \setsansfont{KoHo Custom}[
        Path=\DuneNoirLocalFontPath/,
        UprightFont={KoHo-Regular.ttf},
        BoldFont={KoHo-Bold.ttf},
        ItalicFont={KoHo-Italic.ttf},
        BoldItalicFont={KoHo-BoldItalic.ttf}
      ]
      \setmainfont{KoHo Custom}[
        Path=\DuneNoirLocalFontPath/,
        UprightFont={KoHo-Regular.ttf},
        BoldFont={KoHo-Bold.ttf},
        ItalicFont={KoHo-Italic.ttf},
        BoldItalicFont={KoHo-BoldItalic.ttf}
      ]
    }{
      \DuneNoirSansFallback
    }
    \DuneNoirMonoFallback
  }%
}

\DuneNoirUseHelvetica
\makeatletter
\input{../dune_dark_template_lab/beamerthemeDUNENoir.sty}
\makeatother

\ifPDFTeX
  \usepackage[T1]{fontenc}
  \usepackage[utf8]{inputenc}
\fi

\usepackage{amsmath,amssymb,mathtools}
\usepackage{graphicx,booktabs,array,hyperref}
\usepackage{siunitx}
\usepackage{appendixnumberbeamer}
\usepackage[dvipsnames]{xcolor}

\sisetup{detect-all,per-mode=symbol}
\DeclareSIUnit{\sample}{sample}
\DeclareSIUnit{\tick}{tick}
\DeclareSIUnit{\pe}{PE}

\graphicspath{{../../../daphne_mezz_xc_sim/data/output/plots/}}

\title[XCorr Trigger Timing]{XCorr Trigger Timing}
\subtitle{Relative to raw pulse rising edges}
\author{Manuel Arroyave}
\date{May 12, 2026}
\newcommand{\mydate}{May 12, 2026}

\newcommand{\repo}{\texttt{daphne\_mezz\_xc\_sim}}
\newcommand{\sampleperiod}{\SI{16}{\nano\second}}
\newcommand{\clockfreq}{\SI{62.5}{\mega\hertz}}
\newcommand{\fixedlatency}{\SI{1008}{\nano\second}}
\newcommand{\fixedticks}{\SI{63}{\sample}}
\newcommand{\xcperpe}{\num{581.9}}
\newcommand{\rawbaseline}{\num{2692}}
\newcommand{\rawmax}{\num{16383}}
\newcommand{\rawspan}{\num{13691}}
\newcommand{\rawrangepe}{\num{1369}}

\begin{document}

\begin{frame}
  \titlepage
\end{frame}

\begin{frame}{Main Result}
\small
\begin{block}{Timing decomposition}
\[
\Delta n(P)=n_{\mathrm{trig}}(P)-n_{\mathrm{edge}}
          = n_{\mathrm{latency}} + n_{\mathrm{walk}}(P).
\]
\end{block}

\begin{itemize}
  \item \(n_{\mathrm{edge}}\): raw rising-edge timestamp.
  \item \(n_{\mathrm{trig}}(P)\): xcorr trigger timestamp at PE threshold \(P\).
  \item Fixed xcorr latency: \(63\) samples = \fixedlatency.
  \item Residual walk: \(n_{\mathrm{walk}}(P)\), the threshold-dependent part.
\end{itemize}
\end{frame}

\begin{frame}{Data Sources}
\scriptsize
\begin{columns}[T,onlytextwidth]
\column{0.50\textwidth}
\begin{block}{Raw data}
\begin{itemize}
  \item Real waveform: \path{data/input/run039344_ch35.bin}
  \item Samples: unsigned 14-bit ADC words in 16-bit little-endian storage
  \item Sampling clock: \clockfreq
  \item Sample period: \sampleperiod
\end{itemize}
\end{block}

\column{0.44\textwidth}
\begin{block}{Raw dynamic range}
\[
A_{\max} = \rawmax-\rawbaseline=\rawspan~\mathrm{ADC}
\]
\[
1~\mathrm{PE}=10~\mathrm{ADC}
\quad\Rightarrow\quad
A_{\max}\approx\rawrangepe~\mathrm{PE}
\]
\end{block}
\end{columns}

\vspace{0.2em}
\textbf{Study location:} \path{~/repo/daphne_mezz_xc_sim}
\end{frame}

\begin{frame}{Study Artifacts And PE Calibration}
\footnotesize
\begin{columns}[T,onlytextwidth]
\column{0.56\textwidth}
\begin{block}{Code and outputs}
\begin{itemize}
  \item Timing script: \path{scripts/plot_xcorr_delay_vs_pe.py}
  \item XCorr trace: \path{data/output/analysis/run039344_ch35_rtl_thr2k_len1024_xcorr.txt}
  \item Summary CSV: \path{data/output/analysis/xcorr_threshold_timing/run039344_ch35_delay_vs_pe_summary.csv}
\end{itemize}
\end{block}

\column{0.40\textwidth}
\begin{block}{PE threshold conversion}
The saturated-pulse plateau gives
\[
1~\mathrm{PE}_{\mathrm{thr}} = \xcperpe~\mathrm{xcorr~counts}.
\]
\end{block}

\begin{block}{Plots}
\begin{itemize}
  \item Residual vs PE threshold
  \item Example pulses with edge marks
\end{itemize}
\end{block}
\end{columns}
\end{frame}

\begin{frame}{What Is Compared}
\scriptsize
\begin{columns}[T,onlytextwidth]
\column{0.48\textwidth}
\begin{block}{Raw reference timestamp}
For each pulse:
\[
n_{\mathrm{edge}} = 
\operatorname*{arg\,max}_{n<n_{\mathrm{peak}}}
  \frac{d}{dn}\left(x[n]-\mathrm{median}(x)\right).
\]
This is the peak finder on the rising edge of the raw waveform.
\end{block}

\column{0.48\textwidth}
\begin{block}{XCorr trigger timestamp}
For PE threshold setting \(P\):
\[
T(P)=P\cdot \xcperpe .
\]
\[
n_{\mathrm{trig}}(P)=\min n
\]
with \(y[n]>T\), \(y[n-1]>T\), and \(y[n-2]\le T\). This mirrors the two-cycle threshold-crossing rule used by the emulator.
\end{block}
\end{columns}

\vspace{0.15em}
\textbf{Important alignment point.}
The xcorr timestamp is produced by a causal matched-filter pipeline. It is not in the same phase reference as the raw ADC sample.
\end{frame}

\begin{frame}{Causal Matched Filter}
\scriptsize
\textbf{Filter form}
\[
y[n] = \sum_i h_i\,x[n-d_i],
\qquad
d_i \simeq 3i+4.
\]

\begin{columns}[T,onlytextwidth]
\column{0.48\textwidth}
\begin{block}{Where the delay comes from}
\begin{itemize}
  \item RTL-style transposed FIR tap \(i\) appears at approximately \(3i+4\) samples.
  \item The xcorr phase follows the template centroid, not the raw leading edge.
\end{itemize}
\end{block}

\column{0.48\textwidth}
\begin{block}{Template centroid}
Using the negative template lobe:
\[
\tau_h =
\frac{\sum_{h_i<0}(-h_i)(3i+4)}{\sum_{h_i<0}(-h_i)}
=63.9~\mathrm{samples}.
\]
\[
63~\mathrm{samples}\cdot\sampleperiod=\fixedlatency.
\]
\end{block}
\end{columns}
\end{frame}

\begin{frame}{Subtract The Fixed Latency}
\small
\begin{block}{Latency-corrected residual}
\[
r(P)=\left[n_{\mathrm{trig}}(P)-n_{\mathrm{edge}}\right]-63~\mathrm{samples}.
\]
\[
r_{\mathrm{ns}}(P)=16~r(P).
\]
\end{block}

\vspace{0.3em}
\begin{itemize}
  \item The fixed latency is estimated from the \SI{1}{PE} threshold median: \fixedlatency = \fixedticks.
  \item The residual is the threshold-dependent walk that matters for timing.
  \item The shaded green band in the next plot is \(\pm 2\) ticks = \(\pm\SI{32}{\nano\second}\).
\end{itemize}
\end{frame}

\begin{frame}{Residual Timing Walk Versus PE Threshold}
\centering
\includegraphics[width=\linewidth,height=0.62\textheight,keepaspectratio]{xcorr_residual_vs_pe_run039344_ch35.pdf}

\vspace{0.1em}
\scriptsize\textcolor{noirmuted}{Black line: median residual. Blue band: 10--90\% interval. Green band: \(\pm 2\) ticks.}
\end{frame}

\begin{frame}{Statistical Summary}
\scriptsize
\centering
\setlength{\tabcolsep}{6pt}
\begin{tabular}{rrrrr}
\toprule
\textbf{Threshold} &
\textbf{Matched} &
\textbf{Residual} &
\textbf{Residual p10} &
\textbf{Residual p90} \\
\textbf{(PE)} & \textbf{pulses} & \textbf{(ticks)} & \textbf{(ticks)} & \textbf{(ticks)} \\
\midrule
1    & 386 &  0 & -2 & 15 \\
10   & 384 &  0 & -2 & 15 \\
20   & 382 &  0 & -2 & 16 \\
100  & 364 &  2 &  0 & 17 \\
200  & 331 &  4 &  2 & 17 \\
500  & 263 & 10 &  8 & 14 \\
800  & 230 & 15 & 14 & 19 \\
1000 & 199 & 19 & 18 & 23 \\
1200 & 162 & 25 & 24 & 26 \\
1300 & 138 & 29 & 28 & 29 \\
\bottomrule
\end{tabular}

\vspace{0.2em}
\scriptsize
\textbf{Interpretation:} low-threshold median residual is \(\leq 2\) ticks; large PE thresholds walk later because the trigger waits for a higher point on the xcorr lobe.
\end{frame}

\begin{frame}{Example Pulses: First Pair}
\centering
\includegraphics[
  width=\linewidth,
  height=0.68\textheight,
  keepaspectratio,
  trim=0 389pt 0 0,
  clip
]{xcorr_example_pulses_run039344_ch35.pdf}
\end{frame}

\begin{frame}{Example Pulses: Second Pair}
\centering
\includegraphics[
  width=\linewidth,
  height=0.68\textheight,
  keepaspectratio,
  trim=0 0 0 389pt,
  clip
]{xcorr_example_pulses_run039344_ch35.pdf}
\end{frame}

\begin{frame}{Why Threshold Walk Appears}
\scriptsize
\textbf{Crossing equation.}
Near a pulse, after latency alignment:
\[
y_{\mathrm{aligned}}[k] \approx A\,g[k],
\]
where \(A\) is pulse amplitude and \(g[k]\) is the normalized xcorr lobe.
\[
k_T = \min\{k: A\,g[k] > T\}.
\]

\vspace{0.05em}
\begin{itemize}
  \item If \(T\) is small, \(k_T\) is close to the leading part of the xcorr lobe.
  \item If \(T\) is large, \(k_T\) moves up the lobe toward the peak.
  \item For saturated pulses, raising \(T\) from low PE to \(\sim 1000\) PE moves the median crossing by about \(19\) ticks after latency correction.
\end{itemize}
\end{frame}

\begin{frame}{Practical Consequence}
\small
\begin{alertblock}{Comparator behavior}
The timing comparator must either subtract the fixed latency, use a constant-fraction-like timing reference, or constrain threshold settings to the region where residual walk is acceptable.
\end{alertblock}

\vspace{0.5em}
\begin{itemize}
  \item The 63-sample term is a phase/latency offset.
  \item The residual walk is the physics/algorithm behavior that changes with threshold.
  \item The PE scan identifies the threshold boundaries where that residual leaves the acceptable timing band.
\end{itemize}
\end{frame}

\begin{frame}{Reproducibility Checklist}
\scriptsize
\textbf{Commands used for the PE scan}
\begin{enumerate}
  \item Run or reuse the xcorr trace from \repo.
  \item Execute:
  \[
  \texttt{python3 scripts/plot\_xcorr\_delay\_vs\_pe.py}
  \]
  \item Use the generated summary CSV, residual plot, and pulse examples.
\end{enumerate}

\vspace{0.25em}
\textbf{Validation anchors}
\begin{itemize}
  \item \(\clockfreq \Rightarrow \sampleperiod\) per sample.
  \item Raw dynamic range: \(\rawspan~\mathrm{ADC}\approx\rawrangepe~\mathrm{PE}\).
  \item Filter/template centroid: \(63.9\) samples.
  \item Measured fixed latency from data: \(63\) samples.
\end{itemize}
\end{frame}

\begin{frame}{Generated Files}
\tiny
\textbf{Analysis}
\begin{itemize}
  \item \path{data/output/analysis/xcorr_threshold_timing/run039344_ch35_delay_vs_pe_summary.csv}
  \item \path{data/output/analysis/xcorr_threshold_timing/run039344_ch35_pe_calibration.csv}
\end{itemize}

\vspace{0.35em}
\textbf{Figures}
\begin{itemize}
  \item \path{data/output/plots/xcorr_residual_vs_pe_run039344_ch35.pdf}
  \item \path{data/output/plots/xcorr_example_pulses_run039344_ch35.pdf}
\end{itemize}
\end{frame}

\begin{frame}{Bottom Line}
\small
\begin{enumerate}
  \item The measured xcorr time is fixed matched-filter latency plus residual threshold walk.
  \item The fixed latency is \(\approx 63\) samples, confirmed mathematically and from the \SI{1}{PE} data point.
  \item The timing quantity of interest is the residual:
  \[
    n_{\mathrm{trig}} - n_{\mathrm{edge}} - 63.
  \]
  \item Low thresholds have near-zero median residual.
  \item High thresholds introduce deterministic threshold walk because the crossing point moves later on the xcorr lobe.
\end{enumerate}
\end{frame}

\end{document}
